\documentclass[10pt,a4paper]{amsart}
\usepackage[margin=19mm]{geometry}
\usepackage{amsmath,amssymb,mathtools}
\usepackage[scr=rsfs,cal=euler]{mathalfa}
\usepackage{xfrac}
\usepackage[unicode,hidelinks,bookmarks=false]{hyperref}
\newcommand{\ie}{i.e.\ }
\newcommand{\cf}{cf.\ }
\newcommand{\resp}{respectively\ }
\newcommand{\Subsection}{\subsection}
\providecommand{\nobreakdash}{\nobreak\hbox{-}}
\setlength{\emergencystretch}{2em}
\multlinegap0pt
\usepackage{amssymb}
\usepackage{xcolor}
\usepackage{float}
\usepackage{tikz}
\usepackage{mathtools}
\usepackage[shortlabels]{enumitem}
\newtheorem{theorem}{Theorem}
\newtheorem{lemma}{Lemma}
\renewcommand{\L}{\mathcal{L}}
\newcommand{\weight}{b}
\title{Discrete Laplace and transition operators over non-Archimedean ordered fields}
\author{Anna Muranova}
\date{Source: arXiv:2207.14018v3 (CC BY 4.0)}
\begin{document}
\maketitle

\noindent\emph{Source and licence.} Discrete Laplace and transition operators over non-Archimedean ordered fields. Version arXiv:2207.14018v3 (CC BY 4.0), distributed by arXiv under the Creative Commons Attribution 4.0 International licence, http://creativecommons.org/licenses/by/4.0/, as recorded in the arXiv licence metadata for this exact version and shown on the abstract page https://arxiv.org/abs/2207.14018v3. Full licence notice: \texttt{licenses/arxiv-2207.14018-CC-BY-4.0.txt}.\par\smallskip

\noindent\textit{Editorial scope.} Counted unit: the article's theorem Theorem::Cheeger (source0.tex lines 461--465). The article's complete original proof of it is delivered with it (source0.tex lines 501--592, one \texttt{proof} environment of 92 lines, which the article places away from the statement). No mathematics is written, repaired or re-derived: the statement and the proof are the article's own lines, in the article's own notation, and no retained line is substituted or re-flowed.  Retained uncounted as the source-local context the proof needs: lines 469--482.  Nothing was omitted from any retained span, so the package carries no omission record.  No retained line contains a citation, so the wrapper carries no bibliography and no bibliography entry.  No retained line contains a figure environment, an image-inclusion command or drawing code, so no image is delivered and the package declares no asset dependency.  Editorial formatting only, in addition: an AMS \texttt{amsart} preamble that restates the article's own declarations and macros used by the retained lines, namely the environments \texttt{theorem}, \texttt{lemma} on separate counters, together with the standard packages those lines need.  The retained environments print in this document as Theorem 1, Lemma 1 in order of appearance, whereas the article prints the same items with the numbering of its own full document.  The complete native source of the article as distributed in the arXiv e-print for this exact version is bundled unchanged as \texttt{sources/128816/source0.tex}.
\begin{lemma}\label{Lemma:lambda1g}
\begin{equation}\label{lambda1g}
\lambda_1\succeq \dfrac{\frac{1}{2}\sum_{x,y}\weight (x,y)(g(x)-g(y))^2}{\sum_x \weight (x)g(x)^2},
\end{equation}
where
\begin{equation*}
g(x)=
\begin{cases}
v_1(x), \mbox{ for }v_1(x)\succ 0\; (\Leftrightarrow x\in V^+),\\
0, \mbox{ otherwise } (\Leftrightarrow x\in V^-).
\end{cases}
\end{equation*}
Note that $V^-$ is non-empty, since $\langle v_1, 1\rangle=0$.
\end{lemma}

\begin{theorem}[Cheeger's inequality]\label{Theorem::Cheeger}
\begin{equation*}
\lambda_1\succeq 1-\sqrt{1-h^2}
\end{equation*}
\end{theorem}

\begin{proof}[Proof of Theorem \ref{Theorem::Cheeger}]
Within the proof $V^+, V^-$ and $g$ are as in the proof of Lemma \ref{Lemma:lambda1g}.

Without loss of generality, we can assume that $\weight (V^+)\preceq \weight (V^-)$ (otherwise, consider eigenfunction $-v_1$). Moreover, we can enumerate all $n=\#V$ vertices of the graph so, that 
\begin{equation*}
\underbrace{v_1(x_1)\preceq\dots \preceq v_1(x_{n_0})\preceq 0}_{V^-}\prec\underbrace{ v_1(x_{n_0+1})\preceq\dots\preceq v_1(x_{n})}_{V^+}.
\end{equation*}
We introduce some further notations. Let
\begin{equation*}
S_i = \{x_j\;:\; 1\le j \le i\},\; i= {1,\dots, n}
\end{equation*}
and
\begin{equation*}
\beta = \min_{i= {1,\dots, n}} \dfrac{\weight(\partial S_i)}{\min\{\weight(S_i), \weight(V\setminus S_i)\}}.
\end{equation*}
Then
\begin{equation*}
\weight(\partial S_i) = \sum_{\substack{j: 1\le j \le i\\k: i<k\le n}}  \weight (x_j,x_k)
\end{equation*}
\begin{equation*}
\weight(S_i) = \sum_{j: 1\le j \le  i}  \weight (x_j)\mbox{ and } \weight(V\setminus S_i) = \sum_{j: i< j \le n}  \weight (x_j)
\end{equation*}
and we note that $h\preceq\beta$.



Denoting the r.h.s of \eqref{lambda1g} by $W$ and multiplying numerator and denominator by $\sum_{x,y} \weight (x,y)(g(x)+g(y))^2$ we obtain
\begin{equation*}
W = \dfrac{\frac{1}{2}\sum_{x,y}\weight (x,y)(g(x)-g(y))^2 \sum_{x,y}\weight (x,y) (g(x)+g(y))^2}{\sum_x \weight (x)g(x)^2 \sum_{x,y} \weight (x,y)(g(x)+g(y))^2}
\end{equation*}
and by Cauchy-Schwarz inequality \footnote{The Cauchy-Schwarz inequality in an ordered field can be proven by induction or through Pythagorean theorem} 
\begin{equation}\label{W1}
W\succeq \dfrac{\frac{1}{2}\left(\sum_{x,y}\weight (x,y)\left|g^2(x)-g^2(y)\right|\right)^2}{\sum_x \weight (x)g(x)^2 \sum_{x,y} \weight (x,y)(g(x)+g(y))^2}
\end{equation}
Apart from this from \eqref{lambda1g} we have
\begin{align*}
W\sum_x \weight (x)g(x)^2&= \dfrac{1}{2}\sum_{x,y} \weight (x,y) (g(x)-g(y))^2 =[\mbox{ Green formula}]\\
&=\sum_{x}\L g(x)g(x)\weight (x)=\sum_{x} \weight (x)g^2(x)- \sum_{x,y} \weight (x,y)g(y)g(x)\\
&=[\pm\mbox{the first sum}]\\
&=2\sum_{x} \weight (x)g^2(x)-\sum_{x} \weight (x)g^2(x) - \sum_{x,y} \weight (x,y) g(y)g(x)\\
&=2\sum_{x} \weight (x)g^2(x)-\sum_{x,y} \weight (x,y)g^2(x) -\sum_{x,y} \weight (x,y) g(y)g(x)\\
&=2\sum_{x} \weight (x)g^2(x)-\left(\sum_{x,y} \weight (x,y)g(x)(g(x)+g(y))\right)\\
&=2\sum_{x} \weight (x)g^2(x)-\left(\sum_{x,y} \weight (x,y)g(y)(g(x)+g(y))\right)\\
&=2\sum_{x} \weight (x)g^2(x)-\dfrac{1}{2}\left(\sum_{x,y} \weight (x,y)(g(x)+g(y))^2\right),\\
\end{align*}
where in the last two lines we exchanged notations $x$ and $y$, sum the two lines up and divided by $2$ (usually this trick is used in a proofs of Green's formula).

Therefore, from \eqref{W1} follows
\begin{align*}
W&\succeq \dfrac{\frac{1}{2}\left(\sum_{x,y}\weight (x,y)|g^2(x)-g^2(y)|\right)^2}{\sum_x \weight (x)g(x)^2 (4\sum_{x} \weight (x)g^2(x)-2W\sum_x \weight (x)g(x)^2)}\\
&\succeq \dfrac{\frac{1}{2}\left(\sum_{x,y}\weight (x,y)|g^2(x)-g^2(y)|\right)^2}{\left(\sum_x \weight (x)g(x)^2\right)^2 (4-2W)}
\end{align*}
Now let us estimate the numerator $A:=\frac{1}{2}\left(\sum_{x,y}\weight (x,y)|g^2(x)-g^2(y)|\right)^2$. 
\begin{align*}
A&=\dfrac{1}{2} \left(2\sum_{i<j} \weight (x_i,x_j)(g^2(x_j)-g^2(x_i))\right)^2=2 \left(\sum_{i=1}^{n-1}\sum_{j=i+1}^n \weight (x_i,x_j)(g^2(x_j)-g^2(x_i))\right)^2\\
&=2 \left(\sum_{i=1}^{n-1}\sum_{j=i+1}^n \weight (x_i,x_j)(g^2(x_j)-g^2(x_{j-1})+g^2(x_{j-1})\dots-g^2(x_{i+1})+g^2(x_{i+1})-g^2(x_i))\right)^2\\
&=2 \left(\sum_{i=1}^{n-1}\sum_{j=i+1}^n \weight (x_i,x_j)\sum_{k=i}^{j-1}(g^2(x_{k+1})-g^2(x_k))\right)^2\\
&=2 \left(\sum_{i=1}^{n-1}\sum_{j=i+1}^n\sum_{k=i}^{j-1} \weight (x_i,x_j)(g^2(x_{k+1})-g^2(x_k))\right)^2\\
&=[\mbox {change limits of summation for $k$ and $j$}]\\
&=2 \left(\sum_{i=1}^{n-1}\sum_{k=i}^{n-1}\sum_{j=k+1}^{n} \weight (x_i,x_j)(g^2(x_{k+1})-g^2(x_k))\right)^2\\
&=2 \left(\sum_{i=1}^{n-1}\sum_{k=i}^{n-1}(\weight (x_i,x_{k+1})+\dots+\weight (x_i,x_n))(g^2(x_{k+1})-g^2(x_k))\right)^2\\
&=[\mbox {change limits of summation for $k$ and $i$}]\\
&=2 \left(\sum_{k=1}^{n-1}\sum_{i=1}^{k}(\weight (x_i,x_{k+1})+\dots+\weight (x_i,x_n))(g^2(x_{k+1})-g^2(x_k))\right)^2\\
&=2 \left(\sum_{k=1}^{n-1}\underbrace{\sum_{i=1}^{k}\sum_{j=k+1}^{n}\weight (x_i,x_j)}_{=\weight(\partial S_k)}(g^2(x_{k+1})-g^2(x_k))\right)^2\\
&=2 \left(\sum_{k=1}^{n-1}\weight(\partial S_k)(g^2(x_{k+1})-g^2(x_k))\right)^2=2 \left(\sum_{k=n_0}^{n-1}\weight(\partial S_k)(g^2(x_{k+1})-g^2(x_k))\right)^2.\\
\end{align*}
Note that for $k\ge n_0$ we have $\weight(V\setminus S_k)\preceq \weight( S_k)$ due to $ V^-\subset S_k$. Therefore 
$$
\dfrac{\weight(\partial S_k)}{\weight(V\setminus S_k)}\succeq \beta
$$
for such $k$ and we obtain
\begin{align*}
A&\succeq 2 \left(\sum_{k=n_0}^{n-1}\beta \weight(V\setminus S_k)(g^2(x_{k+1})-g^2(x_k))\right)^2\\
&= 2 \left(\beta\sum_{k=n_0}^{n-1}\sum_{j=k+1}^n \weight (x_j)(g^2(x_{k+1})-g^2(x_k))\right)^2\\
&= 2 \left(\beta\sum_{j=n_0+1}^{n}\sum_{k=n_0}^{j-1} \weight (x_j)(g^2(x_{k+1})-g^2(x_k))\right)^2\\
&= 2 \left(\beta\sum_{j=n_0+1}^{n}\weight (x_j)(g^2(x_{j})-g^2(x_{n_0}))\right)^2= 2 \left(\beta\sum_{j=n_0+1}^{n}\weight (x_j)g^2(x_{j})\right)^2,\\
&= 2 \left(\beta\sum_{j=1}^{n}\weight (x_j)g^2(x_{j})\right)^2,\\
\end{align*}
where the last two equalities are true, since $g(x_i)=0$ for $i\le n_0$ and $V^-\ne \emptyset$.

Therefore,
\begin{align*}
W&\succeq \dfrac{A}{\left(\sum_x \weight (x)g(x)^2\right)^2 (4-2W)}\succeq \dfrac{2\beta^2}{4-2W}\\
\end{align*}
From where follows ($W\preceq 2$ as Rayleigh quotient of $g$), that
$$
W^2-2W+\beta^2\preceq 0
$$
and 
$W\succeq 1-\sqrt{1-\beta^2}$ since discriminant is positive due to $\beta\preceq 1=\dfrac{\weight(\partial S_1)}{\weight(S_1)}$.

\end{proof}

\end{document}
